\songtitle{Líneas grises}

\tag*{1996}
\begin{notes}
Based on the short story \emph{Andrea en el casillero} (1996) by Carlos Thompson.
\end{notes}
\begin{style}
colombian gaita paseo, gaita hembra melody, gaita macho countermelody, maraca drive, caja vallenata groove, upright bass ostinato, hand percussion layers, modal chord drone, tape saturation, plate reverb, mono room vocal, breathy baritone, bittersweet resolve, nocturnal warmth, midtempo sway, 96 BPM, call-and-response chorus, organic live room
\end{style}
\begin{lyrics}
\songpart{Verse 1}
El gris del cielo nublado\\
alivia el encierro frío,\\
mas crece un vacío impío\\
sobre el banco del cansado.\\
El verde en el diagramado\\
miente su fiel trayectoria,\\
y afuera la vida en gloria\\
invita a salir, a ser,\\
pero él decidió saber\\
que hay algo en esta memoria.

\songpart{Chorus}
No puede llorar del todo,\\
que la risa le hace guerra,\\
y en la muerte de la tierra\\
brota un brote de algún modo.

Odio no poder odiar\\
sin que el bien lo contrabalance,\\
que en cada grieta un romance\\
le impide el bien derrumbar.

\songpart{Verse 2}
Se toma un respiro inerte,\\
la mente en blanco a propósito,\\
mas vuelve el mismo compromiso\\
que el alma vuelve más fuerte.\\
Y en ese espacio advierte\\
la voz que quisiera oír:
"¿Cómo 'tash?"—debe existir—\\
y el "vale" de Andrea, mas ella\\
en la finca de aquella estrella,\\
y él sin teléfono, a sufrir.

\songpart{Chorus}
No puede llorar del todo,\\
que la risa le hace guerra,\\
y en la muerte de la tierra\\
brota un brote de algún modo.

Odio no poder odiar\\
sin que el bien lo contrabalance,\\
que en cada grieta un romance\\
le impide el bien derrumbar.

\songpart{Verse 3}
Odia el instante que clara\\
la verdad le hace mirar:\\
no puede del todo odiar\\
la vida que lo ampara.\\
Su carrera, aunque amarga\\
en domingos de circuito,\\
le ha dado el mejor delito:\\
amigos, risas, café,\\
y aunque el gris lo tenga en pie,\\
sabe que es su propio rito.

\songpart{Chorus}
No puede llorar del todo,\\
que la risa le hace guerra,\\
y en la muerte de la tierra\\
brota un brote de algún modo.

Odio no poder odiar\\
sin que el bien lo contrabalance,\\
que en cada grieta un romance\\
le impide el bien derrumbar.
\end{lyrics}
